% Figure from MATH 556 notes, section 9. Compile with the notes preamble.
\begin{tikzpicture}[scale=1.7]
  \fill[black!6] (0,0) circle (1.3);
  \draw[black!55, dashed, thin] (0,0) circle (1.3);
  \node[font=\small] at (-0.55,-0.6) {$M$};
  % x_n radially toward p, y_n along a curve toward p
  \foreach \n in {1,...,7} {
    \pgfmathsetmacro{\r}{1.3*(1-0.55^\n)}
    \pgfmathsetmacro{\ang}{25+100*0.55^\n}
    \draw[black!45, dotted] (25:\r) -- (\ang:\r);
    \fill[figB] (25:\r) circle (0.9pt);
    \node[figB, fill=figB, rectangle, inner sep=0pt, minimum size=2.6pt] at (\ang:\r) {};
  }
  % z_n toward q
  \foreach \n in {1,...,9} {
    \pgfmathsetmacro{\r}{1.3*(1-0.6^\n)}
    \pgfmathsetmacro{\ang}{215-30*0.6^\n}
    \fill[figR] (\ang:\r) circle (0.9pt);
  }
  \draw[figB, fill=white] (25:1.3) circle (1.4pt);
  \draw[figR, fill=white] (215:1.3) circle (1.4pt);
  % a point of M as a constant sequence
  \fill (-0.35,0.55) circle (0.9pt) node[left, font=\scriptsize] {$x = [(x, x, \ldots)]$};
  % labels
  \node[figB, font=\scriptsize, anchor=north west, inner sep=1pt] at (25:0.585) {$x_1$};
  \node[figB, font=\scriptsize, anchor=south, inner sep=2pt] at (80:0.585) {$y_1$};
  \node[figB, font=\scriptsize, anchor=west] at (25:1.36) {$p$: $[\{x_n\}] = [\{y_n\}]$};
  \node[figR, font=\scriptsize, anchor=north] at ({215-30*0.6}:0.52) {$z_1$};
  \node[figR, font=\scriptsize, anchor=east] at (215:1.36) {$q$: $[\{z_n\}]$};
\end{tikzpicture}
